\begin{Example}
\label{ex:fourier-series-distribution}
在本例及下一例中, 利用微分为连续运算这一事实. 考虑 Fourier 级数
\begin{equation}\label{eq:cosine-series-distributional-sum}
\cos x+\cos2x+\cdots+\cos nx+\cdots.
\end{equation}
显然, 这一级数在通常意义下不收敛; 例如, 它在 $x=0$ 处发散. 然而, 级数
\begin{equation}\label{eq:sawtooth-sine-series}
\sin x+\dfrac12\sin2x+\dfrac13\sin3x+\cdots
\end{equation}
在 $(0,2\pi)$ 上收敛到 $(\pi-x)/2$, 并在每个紧子区间上一致收敛, 且级数 \eqref{eq:sawtooth-sine-series} 的部分和在 $\mathbb{R}$ 上一致有界. (这里不证明这些断言; 而是请读者参考 Fourier 级数或高等微积分的教材, 或 Chapter~\ref{chap:function-spaces} 中关于 Fourier 级数的讨论.) 由此可知, \eqref{eq:sawtooth-sine-series} 在分布意义下收敛到 $(\pi-x)/2$ 的 $2\pi$ 周期延拓; 即在 $2\pi$ 的整数倍处有跳跃的"锯齿波". 甚至可以用 Heaviside 函数写出它.\[
\begin{aligned}
&\sin x+\dfrac12\sin2x+\dfrac13\sin3x+\cdots\\
&\quad=\dfrac{\pi-x}2+\pi\sum_{n=1}^\infty H(x-2\pi n)-\pi\sum_{n=0}^\infty H(-2\pi n-x)
\end{aligned}
\]现在通过微分及 dirac delta 的对称性, 得到 \eqref{eq:cosine-series-distributional-sum}:
\begin{equation}\label{eq:cosine-series-dirac-comb}
\begin{aligned}
\cos x+\cos2x+\cos3x+\cdots
&=\dfrac{\mathrm{d}}{\mathrm{d}x}\left(\sin x+\dfrac12\sin2x+\dfrac13\sin3x+\cdots\right)\\
&=-\dfrac12+\pi\sum_{n\in\mathbb{Z}}\delta(x-2\pi n)
\end{aligned}
\end{equation}
(参见 Problem~\ref{ex:piecewise-smooth-distribution-derivative}).
\end{Example}

\begin{Example}
\label{ex:delta-sequence-primitive}
要证明可积函数序列 $f_n:\mathbb{R}\to\mathbb{R}$ 收敛到 delta 函数, 只需证明其原函数收敛到 Heaviside 函数. 以下条件对此是充分的:
\begin{enumerate}\item 对任意 $\epsilon>0$, 有
\begin{equation}\label{eq:delta-sequence-uniform-tail-integrals}
\lim_{n\to\infty}\int_{-\infty}^{-a}f_n(x)\,\mathrm{d}x=0,\quad\lim_{n\to\infty}\int_a^\infty f_n(x)\,\mathrm{d}x=0
\end{equation}
且对 $a\in[\epsilon,\infty)$ 一致成立;
\item 
\begin{equation}\label{eq:delta-sequence-unit-integral-limit}
\lim_{n\to\infty}\int_{-\infty}^\infty f_n(x)\,\mathrm{d}x=1;
\end{equation}

\item $\left|\int_{-\infty}^af_n(x)\,\mathrm{d}x\right|$ 以不依赖于 $a\in\mathbb{R}$ 和 $n\in\mathbb{N}$ 的常数为界.
\end{enumerate}
满足这些条件的函数例子为
\begin{equation}\label{eq:cauchy-kernel-delta-sequence}
f_\epsilon(x)=\dfrac{\epsilon}{\pi(x^2+\epsilon^2)},\quad\epsilon\to0,
\end{equation}

\begin{equation}\label{eq:heat-kernel-delta-sequence}
f_t(x)=\dfrac1{2\sqrt{\pi t}}\exp\left(-\dfrac{x^2}{4t}\right),\quad t\to0^+,
\end{equation}

\begin{equation}\label{eq:sinc-kernel-delta-sequence}
f_n(x)=\dfrac{\sin nx}{\pi x},\quad n\to\infty.
\end{equation}

\end{Example}
